\documentclass[11pt,a4paper]{article}
\usepackage[margin=2.5cm]{geometry}
\usepackage{amsmath,amssymb,amsthm}
\usepackage{booktabs}
\usepackage{listings}
\usepackage[colorlinks=true,linkcolor=blue,citecolor=blue]{hyperref}

\newtheorem{proposition}{Proposition}
\theoremstyle{remark}
\newtheorem{remark}{Remark}

\newcommand{\E}{\mathbb{E}}
\newcommand{\M}{\mathcal{M}}
\newcommand{\one}{\mathbf{1}}
\newcommand{\dd}{\mathrm{d}}

\lstset{basicstyle=\ttfamily\footnotesize,columns=fullflexible,frame=single,
        breaklines=true,showstringspaces=false}

\title{A kicked one-dimensional system:\\
densities, observables and TTCF with jumps}
\author{Working notes on Exercise 9.7.1.1}
\date{7 November 2026}

\begin{document}
\maketitle

\begin{abstract}
We take the one-dimensional system $\dot x = x^2-1$ on $[-1,1]$ of Exercise
9.7.1.1 and add, at every time $T$, an instantaneous ``fire'' $x\mapsto x/K$,
$K>1$. The state space is not restricted: the fire only moves probability
towards $0$, and the support of the density shrinks by itself. We derive the
evolution of the density, three equivalent representations of the expectation
of an observable (direct/Lagrangian, Kawasaki/Eulerian, TTCF), and an exact
transient time correlation function (TTCF) formula that includes the jumps.
All maps are M\"obius transformations, so everything can be checked against
closed-form results. A short Python script (Appendix~\ref{app:code}) verifies
the formulae numerically.
\end{abstract}

\tableofcontents

%======================================================================
\section{The model}
%======================================================================

\subsection{Continuous part}
On the phase space $\M=[-1,1]$ consider
\begin{equation}
  \dot x = G(x) = x^2-1, \qquad \Lambda(x) = G'(x) = 2x .
\end{equation}
The endpoints $\pm1$ are fixed points, $\M$ is invariant forwards and
backwards in time, $-1$ is attracting and $+1$ repelling. Since
$\dot x = -(1-x^2)$, the solution is $x(t)=\tanh(\operatorname{artanh}x_0-t)$,
i.e., with $\tau_t=\tanh t$,
\begin{equation}\label{eq:flow}
  \Phi^t x = \frac{x-\tau_t}{1-\tau_t x}
  = \frac{(1+x)-(1-x)e^{2t}}{(1+x)+(1-x)e^{2t}} ,
\end{equation}
which coincides with Eq.~(9.131) of the book and is valid for all
$x\in\M$, including $x=\pm1$. Its derivative is
\begin{equation}\label{eq:flowjac}
  \frac{\dd\Phi^t x}{\dd x} = \frac{1-\tau_t^2}{(1-\tau_t x)^2}
  = \exp\Big(\int_0^t \Lambda(\Phi^s x)\,\dd s\Big) .
\end{equation}

\subsection{The fire}
At the times $T,2T,3T,\dots$ the state jumps according to
\begin{equation}
  \psi(x) = x/K, \qquad K>1 .
\end{equation}
$\psi$ maps $\M$ \emph{into} $\M$, onto the smaller interval
$\psi(\M)=[-1/K,1/K]$. We do \emph{not} change the phase space: after a fire
the flow~\eqref{eq:flow} continues to act on the whole of $\M$. Physically:
before the fire many trees are likely; after it, states above $1/K$ cannot be
reached, and probability is moved towards $0$ while keeping its shape.

\subsection{Evolution maps}
Let $S_{s,t}$ denote the evolution from time $s$ to time $t\ge s$, with the
convention that a jump at time $kT$ is applied when the trajectory reaches
$kT$; a trajectory that \emph{starts} at $s=kT$ starts just after the $k$-th
fire. We write $S_t=S_{0,t}$. The one-period map and the evolution up to
$t=nT+\tau$, $0\le\tau<T$, are
\begin{equation}
  S = \psi\circ\Phi^T, \qquad S_t = \Phi^\tau\circ S^n .
\end{equation}
The dynamics is $T$-periodic: $S_{s+T,t+T}=S_{s,t}$.

\begin{remark}[M\"obius structure]
A map $x\mapsto (ax+b)/(cx+d)$ is represented by the matrix
$\left(\begin{smallmatrix}a&b\\c&d\end{smallmatrix}\right)$, and composition
corresponds to matrix multiplication. Here
\begin{equation}
  \Phi^t \leftrightarrow \begin{pmatrix}1&-\tau_t\\-\tau_t&1\end{pmatrix},\qquad
  \psi \leftrightarrow \begin{pmatrix}1&0\\0&K\end{pmatrix},\qquad
  S \leftrightarrow \begin{pmatrix}1&-\tau_T\\-K\tau_T&K\end{pmatrix}.
\end{equation}
So $S_t$ is a M\"obius map for every $t$, and its matrix is a product of
$2\times2$ matrices.
\end{remark}

%======================================================================
\section{Densities}
%======================================================================

\subsection{Continuous evolution: Lagrangian and Eulerian forms}
Let $\Lambda_{a,b}(x)=\int_a^b\Lambda(\Phi^s x)\,\dd s$. Under the flow,
\begin{align}
  f_t(x) &= e^{-\Lambda_{-t,0}(x)}\, f_0(\Phi^{-t}x)
  && \text{(Lagrangian, Eq.~(9.141))} \label{eq:lag}\\
         &= e^{\Omega^{(0)}_{-t,0}(x)}\, f_0(x)
  && \text{(Eulerian, Eq.~(9.142))} \label{eq:eul}
\end{align}
where the dissipation function of a density $f$ is
\begin{equation}\label{eq:Omega}
  \Omega^{(f)}(x) = -\Lambda(x) - G(x)\,\frac{f'(x)}{f(x)},
  \qquad f\,\Omega^{(f)} = -\frac{\dd}{\dd x}\big(G f\big),
\end{equation}
and $\Omega^{(0)}=\Omega^{(f_0)}$. Integrating along the trajectory,
\begin{equation}
  \Omega^{(0)}_{-t,0}(x) = -\Lambda_{-t,0}(x) + \ln f_0(\Phi^{-t}x) - \ln f_0(x),
\end{equation}
so \eqref{eq:lag} and \eqref{eq:eul} coincide: in the Eulerian form $f_0$ is
evaluated at $x$, and the back-evolution is hidden inside $\Omega^{(0)}$.

For $f_0=1/2$ one has $\Omega^{(0)}=-\Lambda=-2x$ (Eq.~(9.139)) and
\begin{equation}\label{eq:R}
  f_t(x) = R_t(x)\,f_0(x),\qquad
  R_t(x) = e^{-\Lambda_{-t,0}(x)} = \frac{\dd\Phi^{-t}x}{\dd x}
  = \frac{1-\tau_t^2}{(1+\tau_t x)^2}
  = \frac{4e^{-2t}}{\big[(1+x)+(1-x)e^{-2t}\big]^2} .
\end{equation}

\begin{remark}[A misprint in the book]
In Eq.~(9.137) the factor $2$ of $\Lambda=2x$ is missing: the correct result is
$\int_0^t\Lambda(\Phi^s x)\,\dd s = 2t-2\ln\frac{A+Be^{2t}}{2}$, $A=1+x$,
$B=1-x$. Consequently Eq.~(9.138) should read
$f^{(t)}(x)=\frac12\cdot 4e^{-2t}/(A+Be^{-2t})^2$, which is~\eqref{eq:R}. One
checks that $\int_{-1}^1 R_t(x)\,\tfrac12\,\dd x=1$, whereas the expression
printed in (9.138) integrates to $2te^{-t}/(1-e^{-2t})\neq1$. Also, below
(9.139), ``$f_0$ is constant in $[0,1]$'' should read $[-1,1]$.
\end{remark}

\subsection{The jump}
The density is transported by $\psi$ (Frobenius--Perron):
\begin{equation}\label{eq:jump}
  f_{kT^+}(x) = K\, f_{kT^-}(Kx)\,\one_{\psi(\M)}(x),
  \qquad \psi(\M)=[-1/K,1/K].
\end{equation}
Both the factor $K$ (the Jacobian of $\psi^{-1}$) and the argument $Kx$ are
needed. With them, normalization is automatic,
$\int f_{kT^+}=\int f_{kT^-}=1$, and the shape is preserved: $f_{kT^+}$ is
$f_{kT^-}$ squeezed into $[-1/K,1/K]$. Simply multiplying by $K$ and cutting,
$K f_{kT^-}(x)\one_{\psi(\M)}(x)$, would \emph{not} conserve probability. It
would correspond to removing whole realizations from the ensemble, which is a
different problem (escape, conditioning).

\subsection{Composition}
For any $t$,
\begin{equation}\label{eq:ft}
  f_t(x) = f_0\big(S_t^{-1}x\big)\, J_t(x)\, \one_{S_t(\M)}(x),
  \qquad J_t(x)=\Big|\frac{\dd S_t^{-1}x}{\dd x}\Big| .
\end{equation}
Equivalently, by recursion over one period,
\begin{equation}\label{eq:recursion}
  f_{nT^+}(x) = K\, R_T(Kx)\, f_{(n-1)T^+}\big(\Phi^{-T}(Kx)\big)\,
  \one_{[-1/K,1/K]}(x).
\end{equation}
Each step of~\eqref{eq:recursion} contains exactly one jump, hence one factor
$K$ and the indicator of $\psi(\M)=[-1/K,1/K]$, whatever the number of the
jump. The earlier jumps are hidden in $f_{(n-1)T^+}$. Unrolling the recursion,
with $y_n=x$ and $y_{j-1}=\Phi^{-T}(Ky_j)$,
\begin{equation}\label{eq:unrolled}
  f_{nT^+}(x) = K^n \prod_{j=1}^{n} R_T(Ky_j)\; f_0(y_0)\;
  \prod_{j=1}^{n}\one_{[-1/K,1/K]}(y_j),
  \qquad
  \prod_{j=1}^{n}\one_{[-1/K,1/K]}(y_j) = \one_{S^n(\M)}(x).
\end{equation}
The factor $K^n$ appears explicitly, and the product of indicators is the
indicator of the true support $S^n(\M)$ in~\eqref{eq:ft}. That support is
\emph{not} $[-K^{-n},K^{-n}]$, which would be correct only if the flow were
rescaled after each fire. For example, for $K=2$, $T=0.5$, after two jumps
$S^2(\M)=[-0.3908,\,0.0246]$ (Table~\ref{tab:check}), while
$[-K^{-2},K^{-2}]=[-0.25,0.25]$.

\begin{remark}[Composition of the factors]
The factor from $0$ to $2T$ is \emph{not} $R_T(x)^2$. By the cocycle property,
$R_T(x)\,R_T(\Phi^{-T}x)=R_{2T}(x)$. In the Eulerian language: for the second
interval the initial density is $f_T$, which is not uniform, so the relevant
dissipation function is $\Omega^{(f_T)}\neq\Omega^{(f_0)}$. Multiplying $f_T$
by a constant (such as $K$) does not change $\Omega^{(f_T)}$, because only
$\dd\ln f/\dd x$ enters~\eqref{eq:Omega}.
\end{remark}

\subsection{The support shrinks by itself}
Since $\Phi^T(\M)=\M$, we get $S(\M)=\psi(\M)=[-1/K,1/K]$ and, by
induction, $S^n(\M)=[l_n,u_n]$ with
\begin{equation}
  l_n = S(l_{n-1}),\quad u_n=S(u_{n-1}),\qquad l_0=-1,\ u_0=1 .
\end{equation}
Since $\Phi^T u<u$ for $u\in(-1,1)$, we get $u_n<u_{n-1}/K$, so $u_n<K^{-n}$
for $n\ge2$: the maximum number of trees never recovers. The intervals are
nested, $S^{n}(\M)\subset S^{n-1}(\M)$, and both endpoints converge to the
fixed point $x^*$ of $S$ (Section~\ref{sec:closed}). The shrinking need not be
imposed; it is a consequence of the dynamics.

\paragraph{The support is not $[-K^{-n},K^{-n}]$.}
At each period the endpoints are mapped by the full one-period map
$S=\psi\circ\Phi^T$: first the flow moves them, then the fire divides them by
$K$,
\begin{equation}
  l_n = \Phi^T(l_{n-1})/K,\qquad u_n=\Phi^T(u_{n-1})/K .
\end{equation}
Dividing only by $K$ would give $[-K^{-n},K^{-n}]$. The flow pushes every
interior point down towards $-1$ ($\Phi^T u<u$), so for $n\ge2$
\begin{itemize}
\item the upper edge moves \emph{inside}: $u_n<K^{-n}$;
\item the lower edge moves \emph{outside}: $l_n<-K^{-n}$. These are points
  that the continuous evolution carries out of $[-K^{-n},K^{-n}]$.
\end{itemize}
For $K=2$, $T=0.5$, at $n=2$ the flow first takes $-0.5$ to
$\Phi^{0.5}(-0.5)\approx-0.7815$, and the fire then halves it to $-0.3908$:
\begin{center}
\begin{tabular}{ccc}
\toprule
$n$ & $[l_n,u_n]$ & $[-K^{-n},K^{-n}]$\\
\midrule
1 & $[-0.5000,\ 0.5000]$ & $[-0.5000,\ 0.5000]$\\
2 & $[-0.3908,\ 0.0246]$ & $[-0.2500,\ 0.2500]$\\
3 & $[-0.3612,\ -0.2213]$ & $[-0.1250,\ 0.1250]$\\
4 & $[-0.3528,\ -0.3100]$ & $[-0.0625,\ 0.0625]$\\
5 & $[-0.3503,\ -0.3377]$ & $[-0.0313,\ 0.0313]$\\
\bottomrule
\end{tabular}
\end{center}
The two sequences of intervals also end up in different places:
$[-K^{-n},K^{-n}]$ shrinks to the point $0$, whereas $[l_n,u_n]$ shrinks to
the fixed point $x^*\approx-0.3493$ of $S$. This is why
$\E_{nT^+}[x]\to x^*$, and not to $0$.

%======================================================================
\section{Expectations of observables}
%======================================================================

Let $O$ be any (bounded, measurable) function on $\M$; $O(x)=x$ is the
simplest case. Three exact and equivalent representations are available.

\subsection{Direct (Lagrangian, Heisenberg-like)}
\begin{equation}\label{eq:direct}
  \E_t[O] = \int_\M f_t\,O\,\dd x = \int_\M f_0(z)\, O(S_t z)\,\dd z
  = \E_0[O\circ S_t] .
\end{equation}
Sample initial points from $f_0$, evolve them with flow and jumps, and average
$O$. No exponential factor appears: it is absorbed by the change of variables
$x=S_t z$, whose Jacobian is exactly $J_t$.

\subsection{Kawasaki (Eulerian): the ``new observable'' picture}
Define, for $x\in S_t(\M)$,
\begin{equation}
  \Omega_t(x) = \ln\frac{f_0(S_t^{-1}x)\,J_t(x)}{f_0(x)} .
\end{equation}
Then $f_t = f_0\,e^{\Omega_t}\,\one_{S_t(\M)}$ and
\begin{equation}\label{eq:kawasaki}
  \E_t[O] = \E_0\big[\,O\; e^{\Omega_t}\;\one_{S_t(\M)}\big].
\end{equation}
$\ln J_t$ is the sum of $-\Lambda$ integrated over all flow segments of the
backward trajectory, plus $n\ln K$ for $n$ jumps. For $f_0=1/2$ at $t=nT^+$:
\begin{equation}
  e^{\Omega_{nT}(x)} = K^n \prod_{j=1}^n R_T(K y_j),\qquad
  y_n = x,\quad y_{j-1}=\Phi^{-T}(K y_j).
\end{equation}
This is the corrected version of the idea ``new observable
$O_n = O\times(\text{exponential factor})\times K^n$, same density $f_0$''.
The factor $K^n$ is indeed there. The exponential factor is $R_T$ evaluated at
the back-mapped points, not $R_T(x)^n$, and the indicator of the support is
essential.

\begin{remark}[The role of the boundaries]\label{rem:boundaries}
The Kawasaki form~\eqref{eq:kawasaki} is the change of variables $x=S_tz$ in
the direct average~\eqref{eq:direct}. It works cleanly here because the
boundaries do not leak: $G(\pm1)=0$, so $\pm1$ are fixed points and no
trajectory reaches them in finite time, and the fire maps $\M$ into itself.
Hence
\begin{itemize}
\item $S_t$ maps $\M$ into $\M$ without any extra rule; it is invertible
  from $\M$ onto $S_t(\M)$, and the back-mapped point $S_t^{-1}x$ entering
  $e^{\Omega_t}$ always lies in $\M$;
\item probability is conserved: $\E_t[1]=1$.
\end{itemize}
If the boundaries leaked, one would first have to specify what happens there:
\begin{itemize}
\item \emph{absorbing boundary}: \eqref{eq:kawasaki} survives formally, with
  $\one_{S_t(\M)}$ meaning ``came from $\M$ and never left'', but
  $\E_t[1]<1$ (escape setting);
\item \emph{reflecting boundary or reinjection}: $S_t$ is no longer a smooth
  invertible map, $J_t$ acquires extra contributions, and~\eqref{eq:kawasaki}
  must be rederived.
\end{itemize}
The same issue appears in the response formulae of the next section, through
the boundary term $\big[f_0\,G\,(O\circ\Phi^t)\big]_{-1}^{1}$
in~\eqref{eq:autonomous}: the flux of the initial probability through
$\pm1$, weighted by $O$ at time $t$. That term is exact as written if nothing
special happens at $\pm1$, i.e.\ the flow is defined beyond $\M$ and $\M$ is
only the support of $f_0$. If the dynamics is confined to $\M$ by a boundary
rule, the identity $\frac{\dd}{\dd t}(O\circ\Phi^t)=G\,(O\circ\Phi^t)'$ fails
at the boundary, and the formula must be rederived with that rule. The
boundary term vanishes, and no rule is needed, if $G(\pm1)=0$ (our case) or
$f_0(\pm1)=0$. This is also why the jump term in~\eqref{eq:ttcf} is written in
finite form (Remark~\ref{rem:kick}): the kick flow $g(x)=-x$ has
$g(\pm1)\neq0$.
\end{remark}

\subsection{TTCF with jumps}
This is derived in the next section; see Eq.~\eqref{eq:ttcf}.

%======================================================================
\section{TTCF with jumps}
%======================================================================

\subsection{Reminder: the autonomous case}
For the flow alone,
$\frac{\dd}{\dd t}\E_0[O\circ\Phi^t]=\E_0[G\,(O\circ\Phi^t)']$ because
$\frac{\dd}{\dd t}(O\circ\Phi^t) = G\,(O\circ\Phi^t)'$ (the generator commutes
with the flow). Integrating by parts with~\eqref{eq:Omega},
\begin{equation}\label{eq:autonomous}
  \frac{\dd}{\dd t}\E_t[O] = \E_0\big[(O\circ\Phi^t)\,\Omega^{(0)}\big]
  + \Big[f_0\,G\,(O\circ\Phi^t)\Big]_{-1}^{1}.
\end{equation}
The boundary term vanishes because $G(\pm1)=0$. Integration gives Eq.~(9.144):
\begin{equation}
  \E_t[O]=\E_0[O]+\int_0^t\dd s\;\E_0\big[(O\circ\Phi^s)\,\Omega^{(0)}\big].
\end{equation}
With jumps the argument fails: $S_t=\Phi^\tau\circ S^n$ is not a semigroup in
$t$, and the time derivative cannot be moved to time $0$.

\subsection{Derivation with jumps}
Fix $t$ and consider the expectation at time $t$ of a system prepared in $f_0$
at a \emph{variable} starting time $s\in[0,t]$:
\begin{equation}
  F(s) = \E_0\big[O\circ S_{s,t}\big] = \int_\M f_0(x)\,O(S_{s,t}x)\,\dd x .
\end{equation}
Then $F(0)=\E_t[O]$ and $F(t)=\E_0[O]$.
\begin{itemize}
\item \emph{Between jumps} ($s\neq kT$). The backward equation gives
  $\partial_s (O\circ S_{s,t})(x) = -G(x)\,(O\circ S_{s,t})'(x)$. Integrating
  by parts, the boundary term vanishes again since $G(\pm1)=0$, and
  \begin{equation}
    -F'(s) = \E_0\big[\Omega^{(0)}\;(O\circ S_{s,t})\big].
  \end{equation}
\item \emph{At a jump} $s=kT\le t$. $F$ is discontinuous, since
  $S_{kT^-,t}=S_{kT^+,t}\circ\psi$:
  \begin{equation}
    F(kT^-)-F(kT^+) = \E_0\big[O\circ S_{kT,t}\circ\psi - O\circ S_{kT,t}\big].
  \end{equation}
\end{itemize}
Writing $F(0)-F(t)=-\int F'\,\dd s+\sum_k\big[F(kT^-)-F(kT^+)\big]$, we obtain:

\begin{proposition}[TTCF with jumps]
For $nT\le t<(n+1)T$,
\begin{equation}\label{eq:ttcf}
\boxed{\;
  \E_t[O] = \E_0[O]
  + \int_0^t \dd s\;\E_0\big[\Omega^{(0)}\,(O\circ S_{s,t})\big]
  + \sum_{k=1}^{n} \E_0\big[O\circ S_{kT,t}\circ\psi - O\circ S_{kT,t}\big]
\;}
\end{equation}
where $\Omega^{(0)}=-\Lambda-Gf_0'/f_0$ is the dissipation function of the
\emph{continuous} dynamics with respect to $f_0$, the $s$-integral runs over
the flow intervals (the jump instants have zero measure), and $S_{s,t}$
includes all jumps between $s$ and $t$.
\end{proposition}

For $K=1$ (no jumps) the sum vanishes, $S_{s,t}=\Phi^{t-s}$, and the change of
variables $u=t-s$ gives back Eq.~(9.144).

\paragraph{Interpretation.}
\begin{itemize}
\item The first correction is a TTCF in the usual sense: $\Omega^{(0)}$ is
  evaluated at the initial point, drawn from $f_0$, and correlated with $O$
  at the end of a trajectory that \emph{starts at time $s$}. Since the
  dynamics is not autonomous, the start time matters, and it is no longer a
  single correlation function of the lag $t-s$.
\item Each jump contributes a finite ``jump response'': the difference of
  $O$ at time $t$ between a trajectory that starts at $kT$ from $\psi(x)$ and
  one that starts at $kT$ from $x$, averaged over $x\sim f_0$. It requires no
  dissipation function.
\item Only $f_0$ and $\Omega^{(0)}$ ever appear. One never needs the
  dissipation function of a density with restricted support, and points
  whose backward trajectory leaves $S_t(\M)$ cause no problems.
\end{itemize}

\begin{remark}[Jump as a fast flow; boundary terms]\label{rem:kick}
$\psi$ is the time-$\ln K$ map of the flow $\dot x=g(x)=-x$. Writing the jump
term as $\int_0^{\ln K}\dd\sigma\,\frac{\dd}{\dd\sigma}\E_0[A\circ\psi_\sigma]$,
with $A=O\circ S_{kT,t}$ and $\psi_\sigma x=xe^{-\sigma}$, the same integration
by parts gives a ``kick dissipation function''
$\Omega_{\rm kick}=-(gf_0)'/f_0=1+xf_0'/f_0$. However, now $g(\pm1)\neq0$, so
\begin{equation}
  \frac{\dd}{\dd\sigma}\E_0[A\circ\psi_\sigma]
  = \E_0\big[\Omega_{\rm kick}\,(A\circ\psi_\sigma)\big]
  - f_0(1)A(\psi_\sigma 1) - f_0(-1)A(\psi_\sigma(-1)) .
\end{equation}
The boundary terms vanish only if $f_0(\pm1)=0$, e.g.\ for
$f_0=\frac34(1-x^2)$, but not for $f_0=\frac12$. The finite form used
in~\eqref{eq:ttcf} avoids this subtlety altogether.
\end{remark}

\subsection{Using periodicity: correlation functions to compute}
Since $S_{s+T,t+T}=S_{s,t}$, write $s=mT+\sigma$ with $\sigma\in(0,T)$, and
define
\begin{align}
  C(\sigma,u) &= \E_0\big[\Omega^{(0)}(x)\;O(S_{\sigma,\sigma+u}x)\big],
  && \sigma\in(0,T),\ u\ge0, \label{eq:C}\\
  D(u) &= \E_0\big[O(S_{0,u}(\psi x)) - O(S_{0,u}x)\big], && u\ge0,
  \label{eq:D}
\end{align}
where in $D$ the trajectories start just after a jump (phase $0^+$, next jump
at $T$). Then, for $nT\le t<(n+1)T$,
\begin{equation}\label{eq:ttcf-periodic}
  \E_t[O] = \E_0[O]
  + \sum_{m=0}^{n}\int_0^{\min(T,\,t-mT)}\!\!\dd\sigma\;
      C(\sigma,\,t-mT-\sigma)
  + \sum_{k=1}^{n} D(t-kT).
\end{equation}

\subsection{Implementation}
\begin{enumerate}
\item \textbf{Initial ensemble.} Draw $N$ points $x_i\sim f_0$ (or use
  quadrature nodes and weights). Compute $\Omega^{(0)}(x_i)$ once.
\item \textbf{Phase grid.} Choose $\sigma_1,\dots,\sigma_P\in(0,T)$ (e.g.\
  Gauss--Legendre nodes).
\item \textbf{Flow correlations.} For each $\sigma_p$, evolve every $x_i$
  with the kicked dynamics starting at phase $\sigma_p$ (first jump after a
  time $T-\sigma_p$), recording $O$ on a grid of lags $u$. Accumulate
  $C(\sigma_p,u)\approx\frac1N\sum_i\Omega^{(0)}(x_i)\,O(\cdot)$.
\item \textbf{Jump correlations.} Evolve the pairs $x_i$ and $\psi(x_i)$
  from phase $0^+$ and accumulate the difference of $O$ to obtain $D(u)$.
\item \textbf{Assemble} with~\eqref{eq:ttcf-periodic}, using a quadrature in
  $\sigma$.
\end{enumerate}
\emph{Cost.} In the autonomous case one batch of trajectories gives the whole
response. Here one needs one batch per starting phase $\sigma_p$, plus one
batch of pairs for $D$. In exchange, as in standard TTCF, the correlations
start from zero-mean quantities ($\E_0[\Omega^{(0)}]=0$ by~\eqref{eq:Omega}
with $G(\pm1)=0$; $D$ is a difference), which is what gives TTCF its
favourable signal-to-noise ratio when the perturbation is weak (here: $K$
close to $1$).

\paragraph{Pseudo-code.}
\begin{lstlisting}
x[i] ~ f0, i=1..N ;  W[i] = Omega0(x[i])
for p in 1..P:                      # starting phases sigma_p in (0,T)
    y = x ; time_to_jump = T - sigma_p
    for each lag u on the grid:
        advance y by the flow; apply psi whenever a jump time is crossed
        C[p,u] += mean(W * O(y))
y1 = x ; y2 = psi(x)                # pairs started at phase 0+
for each lag u on the grid:
    advance both (flow + jumps at u = T, 2T, ...)
    D[u] = mean(O(y2) - O(y1))
E_t[O] = E0[O] + sum_m int dsigma C(sigma, t-mT-sigma) + sum_k D(t-kT)
\end{lstlisting}

%======================================================================
\section{Closed-form results for this toy model}\label{sec:closed}
%======================================================================

\subsection{\texorpdfstring{Mean of $x$}{Mean of x}}
If $S_t\leftrightarrow\left(\begin{smallmatrix}a&b\\c&d\end{smallmatrix}\right)$
with $c\neq0$, then for $f_0=\frac12$
\begin{equation}\label{eq:exactmean}
  \E_t[x] = \frac12\int_{-1}^1 \frac{az+b}{cz+d}\,\dd z
  = \frac ac + \frac{bc-ad}{2c^2}\,\ln\left|\frac{d+c}{d-c}\right| .
\end{equation}

\subsection{Fixed point and asymptotic state}
With $\tau=\tanh T$, the fixed point of $S$ solves $x-\tau=Kx(1-\tau x)$,
i.e.\ $K\tau x^2+(1-K)x-\tau=0$. The unique root in $\M$ is
\begin{equation}
  x^* = \frac{(K-1)-\sqrt{(K-1)^2+4K\tau^2}}{2K\tau}\ \in(-1/K,0).
\end{equation}
Since
\begin{equation}
  S'(x) = \frac{1-\tau^2}{K(1-\tau x)^2},\qquad
  0<S'(x^*)<\frac{1-\tau^2}{K}<1,
\end{equation}
$x^*$ is attracting. $S$ is increasing, so every orbit converges
monotonically to $x^*$. The density converges to $\delta(x-x^*)$ at the
times $nT^+$, and to $\delta(x-\Phi^\tau x^*)$ at the phase $\tau$: the system
reaches a time-periodic asymptotic state. Unlike the unkicked system, where
$\E_t[x]\to-1$, here $\E_{nT^+}[x]\to x^*$. The phase-space contraction per
period at $x^*$ is $\ln S'(x^*)=\Lambda_{0,T}(x^*)-\ln K<0$: the flow
contributes $\Lambda_{0,T}$, and each fire contributes $-\ln K$.

\subsection{Numerical check}
Table~\ref{tab:check} compares~\eqref{eq:exactmean} with the direct
average~\eqref{eq:direct}, the TTCF formula~\eqref{eq:ttcf}, and the Kawasaki
form~\eqref{eq:kawasaki}, for $K=2$, $T=0.5$, $O(x)=x$, $f_0=\frac12$. The
script also checks~\eqref{eq:ttcf} for $f_0=\frac34(1-x^2)$, where
$\Omega^{(0)}=-4x$; direct and TTCF agree to all printed digits.

\begin{table}[h]
\centering\small
\setlength{\tabcolsep}{4pt}
\begin{tabular}{rcccc l}
\toprule
$t$ & exact & direct & TTCF + jumps & Kawasaki & support $S_t(\M)$\\
\midrule
0.25 & $-0.1652900760$ & $-0.1652900760$ & $-0.1652900760$ & $-0.1652900760$ & $[-1,1]$\\
0.50 & $-0.1613031127$ & $-0.1613031127$ & $-0.1613031127$ & $-0.1613031127$ & $[-0.5000,0.5000]$\\
0.75 & $-0.3749729624$ & $-0.3749729624$ & $-0.3749729624$ & $-0.3749729624$ & $[-0.6636,0.2907]$\\
1.00 & $-0.2780722033$ & $-0.2780722033$ & $-0.2780722033$ & $-0.2780722033$ & $[-0.3908,0.0246]$\\
1.50 & $-0.3265087657$ & $-0.3265087657$ & $-0.3265087657$ & $-0.3265087657$ & $[-0.3612,-0.2213]$\\
2.00 & $-0.3424778729$ & $-0.3424778729$ & $-0.3424778729$ & $-0.3424778729$ & $[-0.3528,-0.3100]$\\
3.00 & $-0.3487418470$ & $-0.3487418470$ & $-0.3487418470$ & $-0.3487418470$ & $[-0.3496,-0.3459]$\\
5.00 & $-0.3493266560$ & $-0.3493266560$ & $-0.3493266560$ & $-0.3493266560$ & $[-0.3493,-0.3493]$\\
\bottomrule
\end{tabular}
\caption{$K=2$, $T=0.5$, $x^*=-0.3493309143$. The jumps occur at
$t=0.5,1,1.5,\dots$; values at those times are taken just after the jump.}
\label{tab:check}
\end{table}

%======================================================================
\section{Remarks and open directions}
%======================================================================
\begin{itemize}
\item \textbf{Irreversibility.} $\psi$ is injective but not onto $\M$, so the
  kicked dynamics is not invertible on $\M$, and there is no involution $i$
  with $iSi=S^{-1}$. Moreover, the support of $f_t$ is strictly smaller than
  that of $f_0$, so the ergodic consistency condition is violated. Standard
  transient fluctuation relations for $\Omega_t$ should therefore not be
  expected. The response formula~\eqref{eq:ttcf} uses neither reversibility
  nor ergodic consistency.
\item \textbf{Other jump maps.} Nothing in the derivation of~\eqref{eq:ttcf}
  uses the specific form of $\psi$: any measurable $\psi:\M\to\M$ applied at
  known times gives the same structure. Random fire intensities ($K$ random)
  can be handled by averaging the jump terms over $K$; jumps at random times
  (e.g.\ Poisson) by averaging over the jump times. These have not been
  derived in detail here.
\item \textbf{Towards a population model.} If $x$ is the population size,
  the continuous part should contain the growth law (with births). The
  constraint ``never back to the initial value'' is a property of the
  combined flow-plus-jump dynamics, as in the toy model, where
  $\sup S^n(\M)<K^{-n}$. A flow with $G(\pm1)\neq0$ at the boundaries of the
  chosen phase space would add boundary terms to the TTCF integration by
  parts, which must be checked case by case.
\item \textbf{Earlier reading, set aside.} Cutting the density and
  multiplying by $K$ (removing realizations) does not conserve probability.
  It belongs to the escape/conditioning setting (survival probabilities), not
  to the present one.
\end{itemize}

\appendix
%======================================================================
\section{Python verification script}\label{app:code}
%======================================================================
Requires only \texttt{numpy}. Run with \texttt{python3 kicked\_ttcf.py}.
\lstinputlisting[language=Python]{kicked_ttcf.py}

\end{document}
